%--------------------------------------------------------------------------------------
% AI Nyilatkozat
%--------------------------------------------------------------------------------------
%----------------------------------------------------------------------------
\chapter*{Nyilatkozat generatív mesterséges intelligencia alkalmazásáról}
%----------------------------------------------------------------------------

\noindent $\Box$ \textbf{Nem használtam} semmilyen generatív MI segédeszközt.

\noindent $\Box$ \textbf{Használtam} generatív MI segédeszközt. Az MI-vel generált tartalmakat ellenőriztem, a generált kimenetek valóságtartalmáról meggyőződtem, az alábbi táblázatban megfelelően jelöltem minden használatot.


\begin{table}[H]
\scriptsize
\setlength{\tabcolsep}{2pt}
\renewcommand{\arraystretch}{1.1}
\noindent
\begin{tabular}{|p{2.7cm}|p{2.6cm}|p{7.4cm}|p{1.5cm}|}\hline
\textbf{Felhasználási módok} 
  & \textbf{Generatív MI eszköz(ök) neve} 
  & \textbf{Érintett részek (fejezet, oldalszám, hivatkozás)} 
  & \textbf{Használat becsült aránya} \\ \hline\hline

Irodalomkutatás 
  & -- 
  & - 
  & 0\% \\ \hline
\textbf{Prompt lényegi része}  
  & \multicolumn{3}{|p{11.5cm}|}{-} \\ \hline\hline

Programkód generálása 
  & \makecell{OpenAI\\ChatGPT}
  & -- 
  & kb.\ 13\% \\ \hline
\textbf{Prompt lényegi része } 
  & \multicolumn{3}{|p{11.5cm}|}{„Írj Python függvényt, amely  LOO hozzájárulását számolja a robusztussági metrika (calculate\_robustness\_score) alapján.”}\\[2pt]
  & \multicolumn{3}{|p{11.5cm}|}{
  „Készíts függvényt, ami beolvassa a megadott séma szerinti elnevezésű \texttt{.keras} fájlokat, és ezeket kliensmodellekként adja vissza további feldolgozáshoz.”}\\[2pt] 
  & \multicolumn{3}{|p{11.5cm}|}{
  „Hozz létre egy, az alábbi példához hasonló függvényt, de ne csak egy \texttt{calc\_rob} függvényt hívjon meg, hanem ciklusban több metrikaszámító függvényt is (pl.\ \texttt{calc\_loss}, \texttt{calc\_fair}, \texttt{calc\_adv}) minden kliensre...”} \\ \hline\hline

Új ötletek, megoldási javaslatok generálása 
  & -
  & -
  & kb.\ 0\% \\ \hline 
  & \multicolumn{3}{|p{11.5cm}|}{-} \\ \hline\hline

Vázlat létrehozása (szövegstruktúra, vázlatpontok) 
  & \makecell{OpenAI\\ChatGPT}
  & \ref{chap: hatter}. fejezet Háttér  
  & kb.\ 5\% \\ \hline 
  & \multicolumn{3}{|p{11.5cm}|}{„Javasolj alfejezetszerkezetet, egy szakdolgozat Háttér fejezetéhez amely a következő témákat tárgyalja...”} \\ \hline


Szövegblokkok létrehozása 
  & \makecell{OpenAI\\ChatGPT}
  & \ref{chap: bev} Bevezetés, \ref{chap: kiert} Kiértékelés, \ref{tab: global_round_diff} táblázat, \ref{tab: rob_adv_spearman} táblázat, \ref{tab: spearman} táblázat
  & kb.\ 10\% \\ \hline
\textbf{Prompt lényegi része}  
  & \multicolumn{3}{|p{11.5cm}|}{„Fogalmazd át a következő bekezdést úgy, hogy erősebben hangsúlyozza a robusztusság fontosságát a federált tanulásban.”}\\[2pt]
  & \multicolumn{3}{|p{11.5cm}|}{
  „Írj egy rövid, captiont a következő táblázathoz...”} \\ \hline\hline

Képek generálása illusztrációs célból 
  & -- 
  & - 
  & 0\% \\ \hline
\textbf{Prompt lényegi része}  
  & \multicolumn{3}{|p{11.5cm}|}{-} \\ \hline\hline

Adatvizualizáció, grafikonok generálása adatpontok alapján 
  & \makecell{OpenAI\\ChatGPT}
  & \ref{chap: kiert} Kiértékelés;
    \ref{fig:all_iid_rob_adv_fair} ábra,
    \ref{fig:all_iid_dist} ábra,
    \ref{fig:all_noniid_dist} ábra,
    \ref{fig:all_iid_avg} ábra,
    \ref{fig:all_datasets_box_4} ábra,
     Melléklet,
    \ref{fig:all_iid_rob_fair_adv_20} ábra
    \ref{fig:all_datasets_box_20_iid} ábra,
    \ref{fig:all_datasets_box_20_Non-IID} ábra, 
  & kb.\ 8\% \\ \hline
\textbf{Prompt lényegi része}  
  & \multicolumn{3}{|p{11.5cm}|}{„Készíts grafikont a amely boxploton ábrázolja a következőket...”} \\ \hline\hline

Prezentáció készítése 
  & --
  & -
  & kb.\ 0\% \\ \hline
\textbf{Prompt lényegi része}  
  & \multicolumn{3}{|p{11.5cm}|}{-} \\ \hline\hline

Egyéb (fordítás, nyelvi stilizálás) 
  & \makecell{OpenAI\\ChatGPT}
  & \ref{chap: abstarct}{Kivonat}, \ref{chap: kapcs_m} Kapcsolodó munkák
  & kb.\ 2\% \\ \hline
\textbf{Prompt lényegi része}  
  & \multicolumn{3}{|p{11.5cm}|}{„Fordítsd le angolra a következő kivonatot.”}\\[2pt]
  & \multicolumn{3}{|p{11.5cm}|}{
  „Foglald össze magyarul a következő bekezdést...”} \\ \hline\hline

\multicolumn{3}{|p{11.5cm}|}{\textbf{Összesített százalékos érték (a feladat érdemi részére nézve):}} 
  & \textbf{kb.\ 33\%} \\ \hline

\multicolumn{4}{|p{14.3cm}|}{\textbf{Összesített érték rövid, szöveges indoklása:} 
A generatív MI-t támogató eszközként használtam:  kódrészletek és vizualizációs sablonok generálására, a dolgozat szerkezetének megtervezésére, valamint szövegblokkok fordítására és stilisztikai finomítására. A kutatási kérdések megfogalmazását, a kísérleti tervet, az implementáció döntő részét, a futtatásokat és az eredmények értelmezését önállóan végeztem; a generatív MI az érdemi munka körülbelül 38\%-ában nyújtott segítséget, főként támogató, nem helyettesítő jelleggel.} \\ \hline

\end{tabular}

\end{table}


